%! Author = nursace
%! Date = 16.09.2026

\documentclass[aspectratio=169,11pt]{beamer}
% Week 2 Lecture — Python enough for data work
% Course: AI Frameworks (ages 16--18; intuition over math)

% Shared Beamer preamble for AI Frameworks
% Audience: college students ages 16--18 (intuition over math)
% Include with: % Shared Beamer preamble for AI Frameworks
% Audience: college students ages 16--18 (intuition over math)
% Include with: % Shared Beamer preamble for AI Frameworks
% Audience: college students ages 16--18 (intuition over math)
% Include with: \input{preamble}

\usetheme{Madrid}
\usecolortheme{default}

% Clean palette: deep slate + teal accent (distinct from Java navy decks)
\definecolor{LectureNavy}{HTML}{1A365D}
\definecolor{LectureAccent}{HTML}{319795}
\definecolor{LectureGray}{HTML}{4A5568}
\definecolor{CodeBg}{HTML}{F7FAFC}
\definecolor{AlertBg}{HTML}{FFF5F5}
\definecolor{TryBg}{HTML}{F0FFF4}
\definecolor{QuizBg}{HTML}{EBF8FF}
\definecolor{AnswerKeyBg}{HTML}{FFFFF0}
\definecolor{AnswerGold}{HTML}{B7791F}
\definecolor{KeywordBlue}{HTML}{2B6CB0}
\definecolor{StringGreen}{HTML}{276749}
\definecolor{CommentGray}{HTML}{718096}

\setbeamercolor{structure}{fg=LectureNavy}
\setbeamercolor{palette primary}{bg=LectureNavy,fg=white}
\setbeamercolor{palette secondary}{bg=LectureAccent,fg=white}
\setbeamercolor{palette tertiary}{bg=LectureGray,fg=white}
% Madrid title box is dark --- title/subtitle must be light for contrast
\setbeamercolor{title}{fg=white,bg=LectureNavy}
\setbeamercolor{subtitle}{fg=white,bg=LectureNavy}
\setbeamercolor{author}{fg=LectureNavy}
\setbeamercolor{date}{fg=LectureGray}
\setbeamercolor{institute}{fg=LectureGray}
\setbeamercolor{frametitle}{bg=LectureNavy,fg=white}
\setbeamercolor{block title}{bg=LectureNavy,fg=white}
\setbeamercolor{block body}{bg=CodeBg,fg=black}
\setbeamercolor{block title alerted}{bg=LectureAccent,fg=white}
\setbeamercolor{block body alerted}{bg=TryBg,fg=black}
\setbeamercolor{block title example}{bg=LectureGray,fg=white}
\setbeamercolor{block body example}{bg=AlertBg,fg=black}

\setbeamertemplate{navigation symbols}{}
\setbeamertemplate{itemize items}[circle]
\setbeamertemplate{enumerate items}[default]

\usepackage[T1]{fontenc}
\usepackage[utf8]{inputenc}
\usepackage{lmodern}
\usepackage{booktabs}
\usepackage{tabularx}
\usepackage{graphicx}
\usepackage{hyperref}
\usepackage{listings}
\usepackage{xcolor}
\usepackage{tikz}
\usetikzlibrary{positioning,arrows.meta,shapes.geometric}

% listings: Python without shell-escape
\lstdefinestyle{pythonstyle}{
    language=Python,
    basicstyle=\ttfamily\small,
    keywordstyle=\color{KeywordBlue}\bfseries,
    commentstyle=\color{CommentGray}\itshape,
    stringstyle=\color{StringGreen},
    numbers=left,
    numberstyle=\tiny\color{CommentGray},
    numbersep=6pt,
    frame=single,
    rulecolor=\color{LectureGray!40},
    backgroundcolor=\color{CodeBg},
    breaklines=true,
    showstringspaces=false,
    tabsize=4,
    captionpos=b,
}
\lstset{style=pythonstyle}

\newcommand{\pycode}[1]{{\ttfamily\small\detokenize{#1}}}
\newcommand{\trythis}[1]{%
    \begin{alertblock}{Try this}
    #1
    \end{alertblock}%
}
\newcommand{\commonmistake}[1]{%
    \begin{exampleblock}{Common mistake}
    #1
    \end{exampleblock}%
}
\newcommand{\remember}[1]{%
    \begin{block}{Remember}
    #1
    \end{block}%
}

% Formative in-class checks: question slide, then a dedicated Answer slide
\newcommand{\quizpause}{%
    \par\vspace{0.5em}
    {\small\textit{Pause --- answer (clicker / raise hand / neighbour) before the next slide.}}%
}
\newcommand{\quizbox}[1]{%
    \setbeamercolor{block title}{bg=KeywordBlue,fg=white}%
    \setbeamercolor{block body}{bg=QuizBg,fg=black}%
    \begin{block}{Quiz}
    #1
    \end{block}%
    \setbeamercolor{block title}{bg=LectureNavy,fg=white}%
    \setbeamercolor{block body}{bg=CodeBg,fg=black}%
}
\newcommand{\answerbox}[1]{%
    \setbeamercolor{block title}{bg=AnswerGold,fg=white}%
    \setbeamercolor{block body}{bg=AnswerKeyBg,fg=black}%
    \begin{block}{Answer}
    #1
    \end{block}%
    \setbeamercolor{block title}{bg=LectureNavy,fg=white}%
    \setbeamercolor{block body}{bg=CodeBg,fg=black}%
}

% Empty author/date placeholders for the lecturer to fill
\author{}
\date{}


\usetheme{Madrid}
\usecolortheme{default}

% Clean palette: deep slate + teal accent (distinct from Java navy decks)
\definecolor{LectureNavy}{HTML}{1A365D}
\definecolor{LectureAccent}{HTML}{319795}
\definecolor{LectureGray}{HTML}{4A5568}
\definecolor{CodeBg}{HTML}{F7FAFC}
\definecolor{AlertBg}{HTML}{FFF5F5}
\definecolor{TryBg}{HTML}{F0FFF4}
\definecolor{QuizBg}{HTML}{EBF8FF}
\definecolor{AnswerKeyBg}{HTML}{FFFFF0}
\definecolor{AnswerGold}{HTML}{B7791F}
\definecolor{KeywordBlue}{HTML}{2B6CB0}
\definecolor{StringGreen}{HTML}{276749}
\definecolor{CommentGray}{HTML}{718096}

\setbeamercolor{structure}{fg=LectureNavy}
\setbeamercolor{palette primary}{bg=LectureNavy,fg=white}
\setbeamercolor{palette secondary}{bg=LectureAccent,fg=white}
\setbeamercolor{palette tertiary}{bg=LectureGray,fg=white}
% Madrid title box is dark --- title/subtitle must be light for contrast
\setbeamercolor{title}{fg=white,bg=LectureNavy}
\setbeamercolor{subtitle}{fg=white,bg=LectureNavy}
\setbeamercolor{author}{fg=LectureNavy}
\setbeamercolor{date}{fg=LectureGray}
\setbeamercolor{institute}{fg=LectureGray}
\setbeamercolor{frametitle}{bg=LectureNavy,fg=white}
\setbeamercolor{block title}{bg=LectureNavy,fg=white}
\setbeamercolor{block body}{bg=CodeBg,fg=black}
\setbeamercolor{block title alerted}{bg=LectureAccent,fg=white}
\setbeamercolor{block body alerted}{bg=TryBg,fg=black}
\setbeamercolor{block title example}{bg=LectureGray,fg=white}
\setbeamercolor{block body example}{bg=AlertBg,fg=black}

\setbeamertemplate{navigation symbols}{}
\setbeamertemplate{itemize items}[circle]
\setbeamertemplate{enumerate items}[default]

\usepackage[T1]{fontenc}
\usepackage[utf8]{inputenc}
\usepackage{lmodern}
\usepackage{booktabs}
\usepackage{tabularx}
\usepackage{graphicx}
\usepackage{hyperref}
\usepackage{listings}
\usepackage{xcolor}
\usepackage{tikz}
\usetikzlibrary{positioning,arrows.meta,shapes.geometric}

% listings: Python without shell-escape
\lstdefinestyle{pythonstyle}{
    language=Python,
    basicstyle=\ttfamily\small,
    keywordstyle=\color{KeywordBlue}\bfseries,
    commentstyle=\color{CommentGray}\itshape,
    stringstyle=\color{StringGreen},
    numbers=left,
    numberstyle=\tiny\color{CommentGray},
    numbersep=6pt,
    frame=single,
    rulecolor=\color{LectureGray!40},
    backgroundcolor=\color{CodeBg},
    breaklines=true,
    showstringspaces=false,
    tabsize=4,
    captionpos=b,
}
\lstset{style=pythonstyle}

\newcommand{\pycode}[1]{{\ttfamily\small\detokenize{#1}}}
\newcommand{\trythis}[1]{%
    \begin{alertblock}{Try this}
    #1
    \end{alertblock}%
}
\newcommand{\commonmistake}[1]{%
    \begin{exampleblock}{Common mistake}
    #1
    \end{exampleblock}%
}
\newcommand{\remember}[1]{%
    \begin{block}{Remember}
    #1
    \end{block}%
}

% Formative in-class checks: question slide, then a dedicated Answer slide
\newcommand{\quizpause}{%
    \par\vspace{0.5em}
    {\small\textit{Pause --- answer (clicker / raise hand / neighbour) before the next slide.}}%
}
\newcommand{\quizbox}[1]{%
    \setbeamercolor{block title}{bg=KeywordBlue,fg=white}%
    \setbeamercolor{block body}{bg=QuizBg,fg=black}%
    \begin{block}{Quiz}
    #1
    \end{block}%
    \setbeamercolor{block title}{bg=LectureNavy,fg=white}%
    \setbeamercolor{block body}{bg=CodeBg,fg=black}%
}
\newcommand{\answerbox}[1]{%
    \setbeamercolor{block title}{bg=AnswerGold,fg=white}%
    \setbeamercolor{block body}{bg=AnswerKeyBg,fg=black}%
    \begin{block}{Answer}
    #1
    \end{block}%
    \setbeamercolor{block title}{bg=LectureNavy,fg=white}%
    \setbeamercolor{block body}{bg=CodeBg,fg=black}%
}

% Empty author/date placeholders for the lecturer to fill
\author{}
\date{}


\usetheme{Madrid}
\usecolortheme{default}

% Clean palette: deep slate + teal accent (distinct from Java navy decks)
\definecolor{LectureNavy}{HTML}{1A365D}
\definecolor{LectureAccent}{HTML}{319795}
\definecolor{LectureGray}{HTML}{4A5568}
\definecolor{CodeBg}{HTML}{F7FAFC}
\definecolor{AlertBg}{HTML}{FFF5F5}
\definecolor{TryBg}{HTML}{F0FFF4}
\definecolor{QuizBg}{HTML}{EBF8FF}
\definecolor{AnswerKeyBg}{HTML}{FFFFF0}
\definecolor{AnswerGold}{HTML}{B7791F}
\definecolor{KeywordBlue}{HTML}{2B6CB0}
\definecolor{StringGreen}{HTML}{276749}
\definecolor{CommentGray}{HTML}{718096}

\setbeamercolor{structure}{fg=LectureNavy}
\setbeamercolor{palette primary}{bg=LectureNavy,fg=white}
\setbeamercolor{palette secondary}{bg=LectureAccent,fg=white}
\setbeamercolor{palette tertiary}{bg=LectureGray,fg=white}
% Madrid title box is dark --- title/subtitle must be light for contrast
\setbeamercolor{title}{fg=white,bg=LectureNavy}
\setbeamercolor{subtitle}{fg=white,bg=LectureNavy}
\setbeamercolor{author}{fg=LectureNavy}
\setbeamercolor{date}{fg=LectureGray}
\setbeamercolor{institute}{fg=LectureGray}
\setbeamercolor{frametitle}{bg=LectureNavy,fg=white}
\setbeamercolor{block title}{bg=LectureNavy,fg=white}
\setbeamercolor{block body}{bg=CodeBg,fg=black}
\setbeamercolor{block title alerted}{bg=LectureAccent,fg=white}
\setbeamercolor{block body alerted}{bg=TryBg,fg=black}
\setbeamercolor{block title example}{bg=LectureGray,fg=white}
\setbeamercolor{block body example}{bg=AlertBg,fg=black}

\setbeamertemplate{navigation symbols}{}
\setbeamertemplate{itemize items}[circle]
\setbeamertemplate{enumerate items}[default]

\usepackage[T1]{fontenc}
\usepackage[utf8]{inputenc}
\usepackage{lmodern}
\usepackage{booktabs}
\usepackage{tabularx}
\usepackage{graphicx}
\usepackage{hyperref}
\usepackage{listings}
\usepackage{xcolor}
\usepackage{tikz}
\usetikzlibrary{positioning,arrows.meta,shapes.geometric}

% listings: Python without shell-escape
\lstdefinestyle{pythonstyle}{
    language=Python,
    basicstyle=\ttfamily\small,
    keywordstyle=\color{KeywordBlue}\bfseries,
    commentstyle=\color{CommentGray}\itshape,
    stringstyle=\color{StringGreen},
    numbers=left,
    numberstyle=\tiny\color{CommentGray},
    numbersep=6pt,
    frame=single,
    rulecolor=\color{LectureGray!40},
    backgroundcolor=\color{CodeBg},
    breaklines=true,
    showstringspaces=false,
    tabsize=4,
    captionpos=b,
}
\lstset{style=pythonstyle}

\newcommand{\pycode}[1]{{\ttfamily\small\detokenize{#1}}}
\newcommand{\trythis}[1]{%
    \begin{alertblock}{Try this}
    #1
    \end{alertblock}%
}
\newcommand{\commonmistake}[1]{%
    \begin{exampleblock}{Common mistake}
    #1
    \end{exampleblock}%
}
\newcommand{\remember}[1]{%
    \begin{block}{Remember}
    #1
    \end{block}%
}

% Formative in-class checks: question slide, then a dedicated Answer slide
\newcommand{\quizpause}{%
    \par\vspace{0.5em}
    {\small\textit{Pause --- answer (clicker / raise hand / neighbour) before the next slide.}}%
}
\newcommand{\quizbox}[1]{%
    \setbeamercolor{block title}{bg=KeywordBlue,fg=white}%
    \setbeamercolor{block body}{bg=QuizBg,fg=black}%
    \begin{block}{Quiz}
    #1
    \end{block}%
    \setbeamercolor{block title}{bg=LectureNavy,fg=white}%
    \setbeamercolor{block body}{bg=CodeBg,fg=black}%
}
\newcommand{\answerbox}[1]{%
    \setbeamercolor{block title}{bg=AnswerGold,fg=white}%
    \setbeamercolor{block body}{bg=AnswerKeyBg,fg=black}%
    \begin{block}{Answer}
    #1
    \end{block}%
    \setbeamercolor{block title}{bg=LectureNavy,fg=white}%
    \setbeamercolor{block body}{bg=CodeBg,fg=black}%
}

% Empty author/date placeholders for the lecturer to fill
\author{}
\date{}


\title{Week 2: Python Enough for Data}
\subtitle{AI Frameworks}

\begin{document}

%----------------------------------------------------------------------
    \begin{frame}
        \titlepage
    \end{frame}

%----------------------------------------------------------------------
    \begin{frame}{Today's learning goals}
        By the end of this lecture you should be able to:
        \begin{enumerate}
            \item Use \textbf{lists} and \textbf{dicts} for small collections of data
            \item Write simple \textbf{loops} and \textbf{functions} for repeated work
            \item Recognise \pycode{None}, basic errors, and a first look at reading files
            \item Explain why we use \textbf{libraries} (and do not reinvent pandas)
        \end{enumerate}
    \end{frame}

    \begin{frame}{Agenda (80 minutes)}
        \begin{enumerate}
            \item Recap from Week 1
            \item Lists (sequences of values)
            \item Dictionaries (labelled values)
            \item Loops and functions
            \item Files, \pycode{None}, and errors (lightly)
            \item Why libraries matter
        \end{enumerate}
        \vspace{0.5em}
        Short quizzes check understanding---not a graded exam.
    \end{frame}

%======================================================================
    \section{Recap}
%======================================================================

    \begin{frame}{Where we are}
        \begin{itemize}
            \item Week 1: AI / ML ideas, supervised learning, tools
            \item Today: enough Python to talk about \textbf{data} before pandas
            \item Next week: pandas DataFrames (tables done properly)
        \end{itemize}
        \vspace{0.5em}
        \remember{We only learn the Python pieces we need for data work---not every language feature.}
    \end{frame}

%======================================================================
    \section{Lists}
%======================================================================

    \begin{frame}{What is a list?}
        A \textbf{list} stores an ordered sequence of values.
        \vspace{0.5em}
        \begin{itemize}
            \item Square brackets: \pycode{[10, 20, 30]}
            \item Indexing starts at \pycode{0}
            \item Can hold numbers, strings, or mixed types (keep it simple for now)
        \end{itemize}
    \end{frame}

    \begin{frame}[fragile]{Creating and indexing}
        \begin{lstlisting}
scores = [80, 90, 70, 100]
print(scores[0])      # 80  (first)
print(scores[-1])     # 100 (last)
print(len(scores))    # 4

scores.append(85)     # add at the end
print(scores)
        \end{lstlisting}
    \end{frame}

    \begin{frame}[fragile]{Slicing (gentle)}
        \begin{lstlisting}
sizes = [50, 65, 80, 95, 110]
print(sizes[1:4])   # [65, 80, 95]  -- from index 1 up to (not including) 4
print(sizes[:2])    # [50, 65]
print(sizes[2:])    # [80, 95, 110]
        \end{lstlisting}
        \vspace{0.3em}
        You do not need every slice trick today---just know lists can give you a ``chunk''.
    \end{frame}

    \begin{frame}{Lists and data thinking}
        \begin{itemize}
            \item One column of a table $\approx$ a list of values
            \item Example: all house sizes, all exam scores
            \item Later, pandas will hold many columns together cleanly
        \end{itemize}
        \vspace{0.4em}
        \commonmistake{Using \pycode{scores[len(scores)]} --- last valid index is \pycode{len(scores) - 1}.}
    \end{frame}

%--- Quiz 1 ---
    \begin{frame}[fragile]{Quiz 1 --- Lists}
        \begin{lstlisting}
xs = [3, 1, 4, 1, 5]
print(xs[2])
        \end{lstlisting}
        \quizbox{What is printed?}
        \begin{enumerate}
            \item \pycode{3}
            \item \pycode{1}
            \item \pycode{4}
            \item \pycode{5}
        \end{enumerate}
        \quizpause
    \end{frame}

    \begin{frame}{Answer --- Quiz 1}
        \answerbox{\textbf{3.} \pycode{4} --- index \pycode{0}$\rightarrow$3, \pycode{1}$\rightarrow$1, \pycode{2}$\rightarrow$4.}
    \end{frame}

%======================================================================
    \section{Dictionaries}
%======================================================================

    \begin{frame}{What is a dictionary?}
        A \textbf{dict} maps \textbf{keys} to \textbf{values} (like a labelled lookup).
        \vspace{0.5em}
        \begin{itemize}
            \item Curly braces: \pycode{\{"size": 70, "rooms": 3\}}
            \item Keys are often strings
            \item Great for one ``row'' of named fields
        \end{itemize}
    \end{frame}

    \begin{frame}[fragile]{Creating and looking up}
        \begin{lstlisting}
house = {"size": 70, "rooms": 3, "city": "yes"}
print(house["size"])     # 70
print(house.get("price", "unknown"))  # safe default

house["price"] = 200     # add or update
print(house.keys())
        \end{lstlisting}
    \end{frame}

    \begin{frame}{List vs dict (when to use which)}
        \begin{center}
            \begin{tabular}{@{}lll@{}}
                \toprule
                & \textbf{List} & \textbf{Dict} \\
                \midrule
                Access by & position (0, 1, 2, \ldots) & name / key \\
                Good for & a column of values & one record with fields \\
                Example & all scores & one student's info \\
                \bottomrule
            \end{tabular}
        \end{center}
        \vspace{0.5em}
        Real tables = many rows; pandas will handle that better than nested lists of dicts.
    \end{frame}

%--- Quiz 2 ---
    \begin{frame}{Quiz 2 --- Dicts}
        \quizbox{You want to store one movie as title, year, and rating with clear names.
        Which structure fits best?}
        \begin{enumerate}
            \item A list of three numbers only
            \item A dictionary with keys like \pycode{"title"}, \pycode{"year"}, \pycode{"rating"}
            \item A single string with no structure
            \item A matplotlib colour
        \end{enumerate}
        \quizpause
    \end{frame}

    \begin{frame}{Answer --- Quiz 2}
        \answerbox{\textbf{2.} A dictionary --- named fields are easy to read and update.}
    \end{frame}

%======================================================================
    \section{Loops and functions}
%======================================================================

    \begin{frame}{Why loops?}
        Repeat the same idea for every item---without copy-paste.
        \vspace{0.4em}
        \begin{itemize}
            \item Sum all scores
            \item Print every name
            \item Count how many values are above 50
        \end{itemize}
    \end{frame}

    \begin{frame}[fragile]{for loops over lists}
        \begin{lstlisting}
scores = [80, 90, 70]
total = 0
for s in scores:
    total = total + s
print(total)           # 240
print(total / len(scores))  # mean
        \end{lstlisting}
    \end{frame}

    \begin{frame}[fragile]{range for counting}
        \begin{lstlisting}
for i in range(3):
    print(i)    # 0 1 2

# Index + value together:
names = ["Ada", "Alan", "Grace"]
for i, name in enumerate(names):
    print(i, name)
        \end{lstlisting}
    \end{frame}

    \begin{frame}[fragile]{Functions: name a recipe}
        \begin{lstlisting}
def mean(values):
    total = 0
    for v in values:
        total = total + v
    return total / len(values)

print(mean([10, 20, 30]))  # 20.0
        \end{lstlisting}
        \vspace{0.3em}
        \begin{itemize}
            \item Inputs = parameters; output = \pycode{return}
            \item Reuse the same logic on many lists
        \end{itemize}
    \end{frame}

    \begin{frame}{Good habits early}
        \begin{itemize}
            \item Meaningful names: \pycode{scores}, not \pycode{x1}
            \item One function = one clear job
            \item Test with a tiny example you can check by hand
        \end{itemize}
        \vspace{0.4em}
        \trythis{In practice: write \pycode{my_min} and \pycode{my_max} without using Python's built-ins first---then compare.}
    \end{frame}

%--- Quiz 3 ---
    \begin{frame}[fragile]{Quiz 3 --- What does this print?}
        \begin{lstlisting}
def double(n):
    return n * 2

xs = [1, 2, 3]
out = []
for x in xs:
    out.append(double(x))
print(out)
        \end{lstlisting}
        \quizbox{Output?}
        \begin{enumerate}
            \item \pycode{[1, 2, 3]}
            \item \pycode{[2, 4, 6]}
            \item \pycode{6}
            \item An error
        \end{enumerate}
        \quizpause
    \end{frame}

    \begin{frame}{Answer --- Quiz 3}
        \answerbox{\textbf{2.} \pycode{[2, 4, 6]} --- each element is doubled and collected in \pycode{out}.}
    \end{frame}

%======================================================================
    \section{Files, None, errors}
%======================================================================

    \begin{frame}{Reading files (idea first)}
        \begin{itemize}
            \item Data often lives in files (CSV = comma-separated values)
            \item Python can open a text file and read lines
            \item Next week pandas will load CSV in one line---today we see the idea
        \end{itemize}
    \end{frame}

    \begin{frame}[fragile]{Tiny file read (awareness)}
        \begin{lstlisting}
# Conceptual demo -- path must exist on your machine
with open("tiny.csv", "r") as f:
    for line in f:
        print(line.strip())
        \end{lstlisting}
        \vspace{0.3em}
        \pycode{with open(...)} closes the file cleanly when finished.
    \end{frame}

    \begin{frame}[fragile]{None means ``no value''}
        \begin{lstlisting}
result = None
if result is None:
    print("No result yet")

# Missing data in the real world often shows up as empty / NaN later in pandas
        \end{lstlisting}
        \vspace{0.3em}
        Use \pycode{is None} (not \pycode{== None}) as the usual style check.
    \end{frame}

    \begin{frame}[fragile]{Errors are information}
        \begin{lstlisting}
scores = [10, 20]
# print(scores[5])   # IndexError: list index out of range

# print(10 / 0)      # ZeroDivisionError
        \end{lstlisting}
        \vspace{0.3em}
        \begin{itemize}
            \item Read the \textbf{error type} and the \textbf{line number}
            \item Fix one problem; re-run
        \end{itemize}
        \vspace{0.3em}
        \commonmistake{Ignoring the traceback and randomly changing unrelated code.}
    \end{frame}

%--- Quiz 4 ---
    \begin{frame}{Quiz 4 --- None and errors (true/false)}
        \quizbox{True or false?\\[0.4em]
        \pycode{IndexError} often means you used a list index that does not exist.}
        \begin{enumerate}
            \item True
            \item False
        \end{enumerate}
        \quizpause
    \end{frame}

    \begin{frame}{Answer --- Quiz 4}
        \answerbox{\textbf{True.} Check \pycode{len(...)} and remember indexing starts at 0.}
    \end{frame}

%======================================================================
    \section{Why libraries}
%======================================================================

    \begin{frame}{Don't reinvent pandas}
        Hand-written loops are great for learning---but:
        \begin{itemize}
            \item Real tables have many columns and messy types
            \item pandas gives \pycode{read_csv}, filters, groupby, \ldots
            \item matplotlib / scikit-learn solve plotting and models
        \end{itemize}
        \vspace{0.5em}
        Your skill: \textbf{ask the right question} and use the right tool.
    \end{frame}

    \begin{frame}{Mental model for the next weeks}
        \begin{enumerate}
            \item Python basics (this week) --- control the computer
            \item pandas --- control \emph{tables}
            \item plots --- see patterns
            \item scikit-learn --- learn patterns from labelled examples
        \end{enumerate}
    \end{frame}

    \begin{frame}{Quiz 5 --- Libraries}
        \quizbox{Why will we use pandas soon instead of only lists of dicts?}
        \begin{enumerate}
            \item Because lists are illegal in Python
            \item Because tables with many rows/columns are easier and safer with a dedicated library
            \item Because pandas replaces the need to think
            \item Because Jupyter cannot run plain Python
        \end{enumerate}
        \quizpause
    \end{frame}

    \begin{frame}{Answer --- Quiz 5}
        \answerbox{\textbf{2.} Dedicated libraries handle real tabular workflows more clearly and reliably.}
    \end{frame}

%======================================================================
    \section{Wrap-up}
%======================================================================

    \begin{frame}{Key takeaways}
        \begin{enumerate}
            \item Lists = ordered values; dicts = named fields
            \item Loops repeat; functions package reusable recipes
            \item \pycode{None} and error messages are part of normal coding
            \item Libraries (pandas next) save us from reinventing table tools
        \end{enumerate}
    \end{frame}

    \begin{frame}{Next up}
        \begin{block}{Practice (today)}
            Warm-ups on lists/strings · mean/min/max functions · optional tiny CSV
        \end{block}
        \vspace{0.5em}
        \begin{block}{Next lecture (Week 3)}
            pandas DataFrames: load CSV, \pycode{head}, \pycode{shape}, select and filter
        \end{block}
    \end{frame}

    \begin{frame}{Questions?}
        \begin{center}
            {\Large Which feels shakiest: lists, dicts, loops, or functions?}
        \end{center}
    \end{frame}

\end{document}
